\chapter{Strong Solutions}

Let $\Omega\subset\mathbb R^n$ and consider
\begin{equation}\tag{9.1}
Lu=a^{ij}(x)D_{ij}u+b^i(x)D_i u+c(x)u.
\end{equation}
A \emph{strong solution} of
\begin{equation}\tag{9.2}
Lu=f
\end{equation}
is a twice weakly differentiable function satisfying the equation almost
everywhere.

\section*{9.1. Maximum Principles for Strong Solutions}

Suppose $L$ is elliptic.  Write
$\mathcal A=[a^{ij}]$, $\mathcal D=\det\mathcal A$, and
$\mathcal D^*=\mathcal D^{1/n}$.  If $\lambda,\Lambda$ are the least and
greatest eigenvalues of $\mathcal A$, then
$0<\lambda\le\mathcal D^*\le\Lambda$.  The coefficient assumptions for the
maximum principle are
\begin{equation}\tag{9.3}
|b|/\mathcal D^*,\ f/\mathcal D^*\in L^n(\Omega),
\qquad c\le0.
\end{equation}

\begin{theorem}[Maximum principle for strong solutions]
\label{gt-ch09-thm-aleksandrov-maximum-principle}
\dcref{ch9:9.1}\group{chapter-9}\source{gilbarg-trudinger:page-0232}
Let $\Omega$ be bounded and let
$u\in C^0(\overline\Omega)\cap W^{2,n}_{\mathrm{loc}}(\Omega)$ satisfy
$Lu\ge f$ under (9.3).  Then
\begin{equation}\tag{9.4}
\sup_\Omega u\le\sup_{\partial\Omega}u^+
+C\|f/\mathcal D^*\|_{L^n(\Omega)},
\end{equation}
where $C$ depends only on $n$, $\operatorname{diam}\Omega$, and
$\|b/\mathcal D^*\|_{L^n(\Omega)}$.  Without continuity on
$\overline\Omega$, the boundary term may be replaced by
$\limsup_{x\to\partial\Omega}u^+(x)$.
\end{theorem}
\begin{proof}
\sketch For a smooth $u$, apply Lemma~\ref{gt-ch09-lem-weighted-contact-set}
with
\[
g(p)=\bigl(|p|^{n/(n-1)}+\mu^{n/(n-1)}\bigr)^{1-n}.
\]
On the positive upper contact set, $Lu\ge f$ bounds
$-a^{ij}D_{ij}u$ by $|b||Du|+|f|$.  H\"older's inequality and the weighted
normal-map estimate give an exponential bound for
$(\sup_\Omega u-\sup_{\partial\Omega}u^+)/\operatorname{diam}\Omega$;
choosing $\mu=\|f/\mathcal D^*\|_n$ yields (9.4).  Approximate a
$W^{2,n}_{\mathrm{loc}}$ solution on interior exhaustions.  First prove the
limit for uniformly elliptic regularizations, then apply
$L_\eta=L+\eta(1+|b|)\Delta$ and let $\eta\downarrow0$ by dominated
convergence.
\[
\frac{\sup_\Omega u-\sup_{\partial\Omega}u^+}
{\operatorname{diam}\Omega}
\le
\left\{
\exp\left[
\frac{2^{n-2}}{n^n\omega_n}
\int_{\Gamma_u^+}\left(1+\frac{|b|^n}{\mathcal D}\right)-1
\right]
\right\}^{1/n}
\|f/\mathcal D^*\|_{L^n(\Gamma_u^+)}.
\tag{9.14}
\]
For the exhaustion approximation $u_m$ on
$\Omega_\varepsilon=\{x:\dist(x,\partial\Omega)>\varepsilon\}$, the limiting
estimate is
\begin{equation}
\tag{9.15}
\sup_{\Omega_\varepsilon}u
\leq\varepsilon+\sup_{\partial\Omega}u^+
+C\|f/\mathcal D^*\|_{L^n(\Omega_\varepsilon)}.
\end{equation}
\end{proof}

For $u\in C^0(\Omega)$, its upper contact set and normal map are
\begin{equation}
\Gamma_u^+
 =\{y\in\Omega:u(x)\le u(y)+p\cdot(x-y)
\text{ for every }x\in\Omega\text{ and some }p\in\mathbb R^n\}.
\tag{9.5}
\end{equation}
\begin{equation}
\chi_u(y)
 =\{p\in\mathbb R^n:u(x)\le u(y)+p\cdot(x-y)
\text{ for every }x\in\Omega\}.
\tag{9.6}
\end{equation}
Thus $\chi_u(y)\ne\varnothing$ exactly on $\Gamma_u^+$.  If $u\in C^1$,
$\chi_u=Du$ there; if $u\in C^2$, then $D^2u\le0$ there.
For the cone $u(x)=a(1-|x-z|/R)$ on $B_R(z)$,
\begin{equation}
\tag{9.7}
\chi_u(y)=
\begin{cases}
-a(y-z)/(R|y-z|),&y\ne z,\\
B_{a/R}(0),&y=z.
\end{cases}
\end{equation}

\begin{lemma}[Contact-set determinant estimate]
\label{gt-ch09-lem-contact-set-determinant-estimate}
\dcref{ch9:9.2}\group{chapter-9}\source{gilbarg-trudinger:page-0233}
If $u\in C^2(\Omega)\cap C^0(\overline\Omega)$ and
$d=\operatorname{diam}\Omega$, then
\begin{equation}\tag{9.8}
\sup_\Omega u\le\sup_{\partial\Omega}u+
\frac d{\omega_n^{1/n}}
\left(\int_{\Gamma_u^+}|\det D^2u|\right)^{1/n}.
\end{equation}
\end{lemma}
\begin{proof}
\sketch Subtract $\sup_{\partial\Omega}u$.  The area formula for the normal
map, justified by the injective perturbation $\chi_u-\varepsilon I$, gives
\begin{equation}\tag{9.9}
|\chi_u(\Omega)|\le\int_{\Gamma_u^+}|\det D^2u|.
\end{equation}
If $u(y)>0$ is its maximum, compare its graph with the cone having vertex
$(y,u(y))$ and base in $B_d(y)$.  The normal image of this cone contains
$B_{u(y)/d}(0)$, which gives (9.8).
\end{proof}

\begin{lemma}[Contact-set estimate for the principal part]
\label{gt-ch09-lem-contact-set-principal-part}
\dcref{ch9:9.3}\group{chapter-9}\source{gilbarg-trudinger:page-0234}
If $u\in C^2(\Omega)\cap C^0(\overline\Omega)$ and
$d=\operatorname{diam}\Omega$, then
\begin{equation}\tag{9.11}
\sup_\Omega u\le\sup_{\partial\Omega}u+
\frac d{n\omega_n^{1/n}}
\left\|\frac{a^{ij}D_{ij}u}{\mathcal D^*}\right\|_{L^n(\Gamma_u^+)}.
\end{equation}
\end{lemma}
\begin{proof}
\sketch On $\Gamma_u^+$, apply
\[
\det A\det B\le\left(\frac{\operatorname{tr}(AB)}n\right)^n
\tag{9.10}
\]
to $A=-D^2u$ and $B=[a^{ij}]$, then use
Lemma~\ref{gt-ch09-lem-contact-set-determinant-estimate}.
\end{proof}

\begin{lemma}[Weighted contact-set estimate]
\label{gt-ch09-lem-weighted-contact-set}
\dcref{ch9:9.4}\group{chapter-9}\source{gilbarg-trudinger:page-0235}
Let $g\ge0$ be locally integrable on $\mathbb R^n$.  For
$u\in C^2(\Omega)\cap C^0(\overline\Omega)$, set
\[
\widetilde M=
\frac{\sup_\Omega u-\sup_{\partial\Omega}u}{d},
\qquad d=\operatorname{diam}\Omega.
\]
Then
\begin{equation}\tag{9.12}
\int_{B_{\widetilde M}(0)}g
\le\int_{\Gamma_u^+}g(Du)|\det D^2u|
\le\int_{\Gamma_u^+}g(Du)
\left(-\frac{a^{ij}D_{ij}u}{n\mathcal D^*}\right)^n.
\end{equation}
\end{lemma}
\begin{proof}
\sketch The weighted area formula gives
\[
\int_{\chi_u(\Omega)}g
\le\int_{\Gamma_u^+}g(Du)|\det D^2u|.\tag{9.13}
\]
The cone comparison from Lemma~\ref{gt-ch09-lem-contact-set-determinant-estimate}
puts $B_{\widetilde M}(0)$ in the normal image, and (9.10) gives the second
inequality.
\end{proof}

\begin{theorem}[Uniqueness for strong Dirichlet solutions]
\label{gt-ch09-thm-strong-dirichlet-uniqueness}
\dcref{ch9:9.5}\group{chapter-9}\source{gilbarg-trudinger:page-0237}
Let $L$ be elliptic in a bounded domain $\Omega$, with
$|b|/\mathcal D^*\in L^n(\Omega)$ and $c\le0$.  If
$u,v\in W^{2,n}_{\mathrm{loc}}(\Omega)\cap C^0(\overline\Omega)$ satisfy
$Lu=Lv$ in $\Omega$ and $u=v$ on $\partial\Omega$, then $u=v$ in $\Omega$.
\end{theorem}
\begin{proof}
\sketch Apply Theorem~\ref{gt-ch09-thm-aleksandrov-maximum-principle} to
$u-v$ and $v-u$ with zero inhomogeneous term.
\end{proof}

\begin{theorem}[Strong maximum principle for strong solutions]
\label{gt-ch09-thm-strong-maximum-principle}
\dcref{ch9:9.6}\group{chapter-9}\source{gilbarg-trudinger:page-0237}
Assume $L$ is uniformly elliptic and $|b|/\lambda$, $c/\lambda$ are bounded.
If $u\in W^{2,n}_{\mathrm{loc}}(\Omega)$ satisfies $Lu\ge0$, then:
\begin{enumerate}
\item[(i)] if $c=0$, an interior maximum forces $u$ to be constant;
\item[(ii)] if $c\le0$, an interior nonnegative maximum forces $u$ to be
constant.
\end{enumerate}
\end{theorem}
\begin{proof}
\sketch If a nonconstant $u$ attained the prohibited maximum $M$, choose an
annulus joining a region where $u<M$ to a maximizing point.  Add a small
multiple of the standard exponential annular barrier to $M-u$.
Theorem~\ref{gt-ch09-thm-aleksandrov-maximum-principle} makes the perturbed
function positive throughout the annulus, contradicting its value at the
maximizing point.
\end{proof}

\section*{9.2. $L^p$ Estimates: Preliminary Analysis}

Let $K_0$ be a cube, let $f\ge0$ be integrable, and suppose
$\int_{K_0}f\le t|K_0|$.  Repeatedly bisect every subcube $K$ satisfying
\begin{equation}\tag{9.16}
\int_K f\le t|K|.
\end{equation}
The maximal rejected cubes form a disjoint-interior family $\mathcal S$ with
\begin{equation}\tag{9.17}
t<\frac1{|K|}\int_Kf\le2^nt
\qquad(K\in\mathcal S).
\end{equation}
If $F=\bigcup_{K\in\mathcal S}K$ and $G=K_0\setminus F$, then
\begin{equation}\tag{9.18}
f\le t\quad\text{a.e. on }G.
\end{equation}
Writing $\widetilde K$ for the parent of $K$ and
$\widetilde F=\bigcup_{K\in\mathcal S}\widetilde K$, one also has
\begin{equation}\tag{9.19}
\int_{\widetilde F}f\le t|\widetilde F|.
\end{equation}
For $f=\mathbf1_\Gamma$, this implies
\begin{equation}\tag{9.20}
|\Gamma|=|\Gamma\cap\widetilde F|\le t|\widetilde F|.
\end{equation}

\section*{9.3. The Marcinkiewicz Interpolation Theorem}

For a measurable $f$, define
\begin{equation}\tag{9.21}
\mu_f(t)=|\{x\in\Omega:|f(x)|>t\}|,
\qquad t>0.
\end{equation}

\begin{lemma}[Distribution-function identities]
\label{gt-ch09-lem-distribution-function-identities}
\dcref{ch9:9.7}\group{chapter-9}\source{gilbarg-trudinger:page-0239}
For $p>0$ and $|f|^p\in L^1(\Omega)$,
\begin{equation}
\mu_f(t)\le t^{-p}\int_\Omega|f|^p.
\tag{9.22}
\end{equation}
\begin{equation}
\int_\Omega|f|^p=p\int_0^\infty t^{p-1}\mu_f(t)\,dt.
\tag{9.23}
\end{equation}
\end{lemma}
\begin{proof}
\sketch The first inequality follows by integrating over $\{|f|>t\}$.
For the second, write
$|f(x)|^p=\int_0^{|f(x)|}pt^{p-1}\,dt$ and apply Fubini.
\end{proof}

\begin{theorem}[Marcinkiewicz interpolation]
\label{gt-ch09-thm-marcinkiewicz-interpolation}
\dcref{ch9:9.8}\group{chapter-9}\source{gilbarg-trudinger:page-0240}
Let $1\le q<r<\infty$, and let $T$ map
$L^q(\Omega)\cap L^r(\Omega)$ linearly into itself.  Suppose
\begin{equation}\tag{9.24}
\mu_{Tf}(t)\le\left(\frac{T_1\|f\|_q}{t}\right)^q,
\qquad
\mu_{Tf}(t)\le\left(\frac{T_2\|f\|_r}{t}\right)^r
\end{equation}
for every $f$ in that intersection and $t>0$.  If $q<p<r$ and
\[
\frac1p=\frac\alpha q+\frac{1-\alpha}r,
\]
then $T$ extends boundedly to $L^p(\Omega)$ and
\begin{equation}\tag{9.25}
\|Tf\|_p\le C T_1^\alpha T_2^{1-\alpha}\|f\|_p,
\end{equation}
where $C=C(p,q,r)$.
\end{theorem}
\begin{proof}
\sketch Split $f=f_1+f_2$ at the level $s$, with $f_1$ supported where
$|f|>s$.  Linearity and (9.24) bound $\mu_{Tf}(t)$ by the weak estimates for
$f_1$ and $f_2$.  Insert this bound into (9.23), choose $t=As$, and use
Fubini to obtain
\[
\|Tf\|_p^p\le C\left(T_1^qA^{p-q}+T_2^rA^{p-r}\right)\|f\|_p^p.
\]
Balancing the two terms in $A$ gives (9.25); density gives the extension.
\end{proof}

\section*{9.4. The Calderon--Zygmund Inequality}

For $f\in L^p(\Omega)$, its Newtonian potential is
\begin{equation}\tag{9.26}
Nf(x)=\int_\Omega\Gamma(x-y)f(y)\,dy.
\end{equation}

\begin{theorem}[Calderon--Zygmund estimate for Newtonian potentials]
\label{gt-ch09-thm-calderon-zygmund-newtonian}
\dcref{ch9:9.9}\group{chapter-9}\source{gilbarg-trudinger:page-0242}
Let $1<p<\infty$ and $w=Nf$ for $f\in L^p(\Omega)$.  Then
$w\in W^{2,p}(\Omega)$, $\Delta w=f$ almost everywhere, and
\begin{equation}\tag{9.27}
\|D^2w\|_p\le C(n,p)\|f\|_p.
\end{equation}
For $p=2$,
\begin{equation}\tag{9.28}
\int_{\mathbb R^n}|D^2w|^2=\int_\Omega f^2.
\end{equation}
For unbounded $\Omega$, the same conclusion holds when $n\ge3$.
\end{theorem}
\begin{proof}
\sketch For smooth compactly supported $f$, integrate by parts on expanding
balls to prove (9.28); approximation handles general $L^2$ data.  For
$Tf=D_{ij}Nf$, (9.28) gives weak type $(2,2)$.  A Calderon--Zygmund cube
decomposition $f=g+b$, with $g$ bounded and each component of $b$ having
mean zero. Its selected cubes $K_l$ satisfy
\begin{equation}
\tag{9.31}
t<\frac1{|K_l|}\int_{K_l}|f|<2^nt,
\end{equation}
and give weak type $(1,1)$:
\[
\mu_{Tf}(t)\le\left(\frac{\|f\|_2}{t}\right)^2.\tag{9.29}
\]
\[
\mu_{Tf}(t)\le C\|f\|_1/t.\tag{9.30}
\]
Theorem~\ref{gt-ch09-thm-marcinkiewicz-interpolation} yields the estimate for
$1<p\le2$:
\begin{equation}
\tag{9.32}
\|Tf\|_p\leq C(n,p)\|f\|_p.
\end{equation}
The symmetry
$\int(Tf)g=\int f(Tg)$ gives $p>2$ by duality, and approximation completes
the proof.
\end{proof}

\begin{corollary}[Calderon--Zygmund estimate for Poisson's equation]
\label{gt-ch09-cor-calderon-zygmund-poisson}
\dcref{ch9:9.10}\group{chapter-9}\source{gilbarg-trudinger:page-0247}
If $\Omega\subset\mathbb R^n$, $u\in W^{2,p}_0(\Omega)$, and
$1<p<\infty$, then
\begin{equation}\tag{9.33}
\|D^2u\|_p\le C(n,p)\|\Delta u\|_p.
\end{equation}
For $p=2$,
\begin{equation}\tag{9.34}
\|D^2u\|_2=\|\Delta u\|_2.
\end{equation}
\end{corollary}
\begin{proof}
\sketch Extend $\Delta u$ by zero and apply
Theorem~\ref{gt-ch09-thm-calderon-zygmund-newtonian} to its Newtonian
potential.  The difference from $u$ is harmonic with zero Sobolev trace, hence
vanishes; density treats general $u$.
\end{proof}

\section*{9.5. $L^p$ Estimates}

For this section assume, for fixed $\lambda,\Lambda>0$,
\begin{equation}\tag{9.35}
a^{ij}\in C^0(\Omega),\quad b^i,c\in L^\infty(\Omega),\quad
a^{ij}\xi_i\xi_j\ge\lambda|\xi|^2,\quad
|a^{ij}|,|b^i|,|c|\le\Lambda.
\end{equation}

\begin{theorem}[Interior $W^{2,p}$ estimate]
\label{gt-ch09-thm-interior-w2p-estimate}
\dcref{ch9:9.11}\group{chapter-9}\source{gilbarg-trudinger:page-0247}
Let $1<p<\infty$ and
$u\in W^{2,p}_{\mathrm{loc}}(\Omega)\cap L^p(\Omega)$ solve $Lu=f$, where
$f\in L^p(\Omega)$ and (9.35) holds.  For every
$\Omega'\Subset\Omega$,
\begin{equation}\tag{9.36}
\|u\|_{2,p;\Omega'}\le
C\bigl(\|u\|_{p;\Omega}+\|f\|_{p;\Omega}\bigr).
\end{equation}
The constant depends on $n,p,\lambda,\Lambda,\Omega',\Omega$ and the moduli
of continuity of $a^{ij}$ on $\Omega'$.
\end{theorem}
\begin{proof}
\sketch Freeze the principal coefficients at $x_0$ and transform the resulting
constant-coefficient operator to the Laplacian.  Corollary
\ref{gt-ch09-cor-calderon-zygmund-poisson} gives
\begin{equation}
\tag{9.37}
\|D^2v\|_{p;\Omega}\leq\frac C\lambda\|L_0v\|_{p;\Omega}.
\end{equation}
Continuity of $a^{ij}$ absorbs the coefficient error on sufficiently small
balls. Applying this to $\eta u$ and setting
$\Phi_k=\sup_{0<\sigma<1}(1-\sigma)^kR^k
\|D^ku\|_{p;B_{\sigma R}}$ gives
\begin{equation}
\tag{9.38}
\Phi_2\leq C(R^2\|f\|_{p;B_R}+\Phi_1+\Phi_0),
\end{equation}
while interpolation gives
\begin{equation}
\tag{9.39}
\Phi_1\leq\varepsilon\Phi_2+C\varepsilon^{-1}\Phi_0.
\end{equation}
Absorb $\Phi_1$ to obtain
\begin{equation}\tag{9.40}
\|D^2u\|_{p;B_{\sigma R}}\le
\frac C{(1-\sigma)^2R^2}
\left(R^2\|f\|_{p;B_R}+\|u\|_{p;B_R}\right).
\end{equation}
A finite covering of $\Omega'$ proves (9.36).
\end{proof}

For a boundary portion $T\subset\partial\Omega$, the condition $u=0$ on
$T$ in the sense of $W^{1,p}(\Omega)$ means that $u$ is the
$W^{1,p}$ limit of $C^1$ functions vanishing near $T$.

\begin{lemma}[Flat-boundary Poisson estimate]
\label{gt-ch09-lem-flat-boundary-poisson}
\dcref{ch9:9.12}\group{chapter-9}\source{gilbarg-trudinger:page-0250}
Let $\Omega^+=\Omega\cap\{x_n>0\}$.  Suppose
$u\in W^{1,1}_0(\Omega^+)$, $f\in L^p(\Omega^+)$, $1<p<\infty$, satisfy
$\Delta u=f$ weakly, and $u=0$ near the curved portion
$(\partial\Omega)^+$.  Then
$u\in W^{2,p}(\Omega^+)\cap W^{1,p}_0(\Omega^+)$ and
\begin{equation}\tag{9.41}
\|D^2u\|_{p;\Omega^+}\le C(n,p)\|f\|_{p;\Omega^+}.
\end{equation}
\end{lemma}
\begin{proof}
\sketch Extend $u$ and $f$ by zero outside $\Omega^+$ and then oddly across
$\{x_n=0\}$.  Testing away from the plane and letting the excluded strip
shrink shows that the extensions satisfy $\Delta u=f$ weakly on
$\mathbb R^n$.  Mollification and Corollary
\ref{gt-ch09-cor-calderon-zygmund-poisson} give (9.41), while oddness gives
the zero trace.
\end{proof}

\begin{theorem}[Local boundary $W^{2,p}$ estimate]
\label{gt-ch09-thm-local-boundary-w2p-estimate}
\dcref{ch9:9.13}\group{chapter-9}\source{gilbarg-trudinger:page-0251}
Let $T\subset\partial\Omega$ be a $C^{1,1}$ boundary portion.  Suppose
$u\in W^{2,p}(\Omega)$, $1<p<\infty$, solves $Lu=f$, vanishes on $T$ in
the $W^{1,p}$ sense, and $L$ satisfies (9.35) with
$a^{ij}\in C^0(\Omega\cup T)$.  For every
$\Omega'\Subset\Omega\cup T$,
\begin{equation}\tag{9.42}
\|u\|_{2,p;\Omega'}\le
C\bigl(\|u\|_{p;\Omega}+\|f\|_{p;\Omega}\bigr).
\end{equation}
The constant depends on $n,p,\lambda,\Lambda,T,\Omega',\Omega$ and the
moduli of continuity of $a^{ij}$ on $\Omega'$.
\end{theorem}
\begin{proof}
\sketch Flatten $T$ by a $C^{1,1}$ diffeomorphism.  The transformed operator
retains (9.35), with controlled constants, and the transformed solution has
zero trace on a flat boundary.  Repeat the frozen-coefficient proof of
Theorem~\ref{gt-ch09-thm-interior-w2p-estimate}, using
Lemma~\ref{gt-ch09-lem-flat-boundary-poisson} on half-balls.  Combine the
boundary estimates with the interior estimate and a finite covering.
\end{proof}

\begin{theorem}[Coercive global $W^{2,p}$ estimate]
\label{gt-ch09-thm-coercive-global-w2p-estimate}
\dcref{ch9:9.14}\group{chapter-9}\source{gilbarg-trudinger:page-0252}
Let $\Omega$ be a $C^{1,1}$ domain and suppose $L$ satisfies (9.35) with
$a^{ij}\in C^0(\overline\Omega)$.  There are $C,\sigma_0>0$, depending only
on $n,p,\lambda,\Lambda,\Omega$ and the moduli of continuity of $a^{ij}$,
such that, for $u\in W^{2,p}(\Omega)\cap W^{1,p}_0(\Omega)$ and
$\sigma\ge\sigma_0$,
\begin{equation}\tag{9.43}
\|u\|_{2,p;\Omega}\le C\|Lu-\sigma u\|_{p;\Omega}.
\end{equation}
\end{theorem}
\begin{proof}
\sketch On $\Omega\times(-1,1)$ use
$L_0=L+\partial_{tt}$ and
$v(x,t)=u(x)\cos(\sqrt\sigma\,t)$.  Theorem
\ref{gt-ch09-thm-local-boundary-w2p-estimate} applied on a thin central
slab controls $\partial_{tt}v$ and yields
\[
\|u\|_{p;\Omega}\le C\|Lu-\sigma u\|_{p;\Omega}\tag{9.44}
\]
for large $\sigma$.  Insert this in the global form of (9.42).
\end{proof}

In Theorems~\ref{gt-ch09-thm-interior-w2p-estimate},
\ref{gt-ch09-thm-local-boundary-w2p-estimate}, and
\ref{gt-ch09-thm-coercive-global-w2p-estimate}, the boundedness assumptions
on the lower-order coefficients may be replaced by
\[
b^i\in L^q,\quad
q>n\text{ if }p\le n,\quad q=p\text{ if }p>n,
\]
and
\[
c\in L^r,\quad
r>n/2\text{ if }p\le n/2,\quad r=p\text{ if }p>n/2.
\]

\section*{9.6. The Dirichlet Problem}

\begin{theorem}[Strong Dirichlet problem in $W^{2,p}$]
\label{gt-ch09-thm-strong-dirichlet-w2p}
\dcref{ch9:9.15}\group{chapter-9}\source{gilbarg-trudinger:page-0253}
Let $\Omega$ be a $C^{1,1}$ domain.  Suppose $L$ is strictly elliptic with
$a^{ij}\in C^0(\overline\Omega)$, $b^i,c\in L^\infty(\Omega)$, and
$c\le0$.  For $1<p<\infty$, $f\in L^p(\Omega)$, and
$\varphi\in W^{2,p}(\Omega)$, there is a unique $u\in W^{2,p}(\Omega)$ such
that
\[
Lu=f\quad\text{in }\Omega,
\qquad u-\varphi\in W^{1,p}_0(\Omega).
\]
\end{theorem}
\begin{proof}
\sketch Subtract $\varphi$.  Approximate the principal coefficients uniformly
and, when $p<2$, approximate $f$ in $L^p$ by $L^2$ data.  The regularity
bootstrap in Lemma~\ref{gt-ch09-lem-boundary-integrability-bootstrap} supplies
the approximating solutions, and Lemma
\ref{gt-ch09-lem-dirichlet-apriori-w2p} bounds them uniformly in $W^{2,p}$.
Weak compactness gives a solution of the limit equation.  Applying the
bootstrap to a homogeneous difference and then
Theorem~\ref{gt-ch09-thm-strong-dirichlet-uniqueness} proves uniqueness.
\end{proof}

\begin{lemma}[Boundary integrability bootstrap]
\label{gt-ch09-lem-boundary-integrability-bootstrap}
\dcref{ch9:9.16}\group{chapter-9}\source{gilbarg-trudinger:page-0253}
Under the hypotheses of
Theorem~\ref{gt-ch09-thm-local-boundary-w2p-estimate}, suppose additionally
$f\in L^q(\Omega)$ for some $q>p$.  Then
$u\in W^{2,q}_{\mathrm{loc}}(\Omega\cup T)$, its trace vanishes on $T$ in
$W^{1,q}$, and (9.42) holds with $p$ replaced by $q$.
\end{lemma}
\begin{proof}
\sketch Localize by $v=\eta u$ and freeze the principal coefficients.
After transforming the frozen operator to $\Delta$, write
\begin{equation}\tag{9.45}
v=Tv+h,
\end{equation}
where the Calderon--Zygmund estimate makes $\|T\|\le1/2$ on a sufficiently
small ball.  The contraction principle improves $W^{2,p}$ to
$W^{2,r}$, where
\[
\frac1r=\max\left\{\frac1q,\frac1p-\frac1n\right\}.
\]
Iteration reaches $q$.  Flattening the boundary and using half-balls gives
the same conclusion near $T$.
\end{proof}

\begin{lemma}[Dirichlet a priori $W^{2,p}$ estimate]
\label{gt-ch09-lem-dirichlet-apriori-w2p}
\dcref{ch9:9.17}\group{chapter-9}\source{gilbarg-trudinger:page-0254}
Under the hypotheses of
Theorem~\ref{gt-ch09-thm-strong-dirichlet-w2p}, there is a constant $C$,
independent of $u$, such that
\begin{equation}\tag{9.46}
\|u\|_{2,p;\Omega}\le C\|Lu\|_{p;\Omega}
\end{equation}
for every $u\in W^{2,p}(\Omega)\cap W^{1,p}_0(\Omega)$.
\end{lemma}
\begin{proof}
\sketch Otherwise normalize a sequence by $\|u_m\|_p=1$ while
$\|Lu_m\|_p\to0$.  The boundary estimate gives a uniform $W^{2,p}$ bound.
Compactness in $L^p$ and weak compactness in $W^{2,p}$ produce a nonzero
limit $u$ with $Lu=0$ and zero trace, contradicting uniqueness.
\end{proof}

\begin{corollary}[Strong Dirichlet problem with continuous boundary data]
\label{gt-ch09-cor-strong-dirichlet-continuous-data}
\dcref{ch9:9.18}\group{chapter-9}\source{gilbarg-trudinger:page-0255}
Under the coefficient and domain assumptions of
Theorem~\ref{gt-ch09-thm-strong-dirichlet-w2p}, let
$p>n/2$, $f\in L^p(\Omega)$, and
$\varphi\in C^0(\partial\Omega)$.  Then
\[
Lu=f\quad\text{in }\Omega,\qquad u=\varphi\quad\text{on }\partial\Omega
\]
has a unique solution
$u\in W^{2,p}_{\mathrm{loc}}(\Omega)\cap C^0(\overline\Omega)$.
\end{corollary}
\begin{proof}
\sketch Approximate $\varphi$ uniformly on the boundary by traces of
$W^{2,p}$ functions and solve using
Theorem~\ref{gt-ch09-thm-strong-dirichlet-w2p}.  The maximum principle makes
the solutions uniformly Cauchy, while Theorem
\ref{gt-ch09-thm-interior-w2p-estimate} and Lemma
\ref{gt-ch09-lem-boundary-integrability-bootstrap} give local $W^{2,p}$
convergence.  The homogeneous uniqueness argument is valid for arbitrary
domains and every $p>1$.
\end{proof}

\begin{theorem}[Higher regularity for strong solutions]
\label{gt-ch09-thm-higher-regularity-strong-solutions}
\dcref{ch9:9.19}\group{chapter-9}\source{gilbarg-trudinger:page-0255}
Let $1<p,q<\infty$, $k\ge1$, and $0<\alpha<1$, and suppose
$u\in W^{2,p}_{\mathrm{loc}}(\Omega)$ solves the elliptic equation $Lu=f$.
The following alternatives hold:
\begin{enumerate}
\item[(i)] if the coefficients of $L$ are in $C^{k-1,1}(\Omega)$ and
$f\in W^{k,q}_{\mathrm{loc}}(\Omega)$, then
$u\in W^{k+2,q}_{\mathrm{loc}}(\Omega)$;
\item[(ii)] if the coefficients and $f$ are in
$C^{k-1,\alpha}(\Omega)$, then
$u\in C^{k+1,\alpha}_{\mathrm{loc}}(\Omega)$.
\end{enumerate}
The conclusions are global if, respectively,
$\partial\Omega\in C^{k+1,1}$ or $C^{k+1,\alpha}$, $L$ is strictly elliptic
with coefficients in the corresponding class on $\overline\Omega$, and
$f\in W^{k,q}(\Omega)$ or $C^{k-1,\alpha}(\overline\Omega)$.
\end{theorem}
\begin{proof}
\sketch Iterate difference-quotient estimates and the local or boundary
$W^{2,q}$ estimate for (i).  For (ii), bootstrap the corresponding Schauder
estimate.  Boundary flattening in the stated regularity class gives the
global alternatives.
\end{proof}

\section*{9.7. A Local Maximum Principle}

For the remainder of the chapter assume $L$ is strictly elliptic with bounded
coefficients and fix $\gamma,\nu$ such that
\begin{equation}\tag{9.47}
\Lambda/\lambda\le\gamma,
\qquad (|b|/\lambda)^2,\ |c|/\lambda\le\nu.
\end{equation}

\begin{theorem}[Local maximum principle]
\label{gt-ch09-thm-local-maximum-principle}
\dcref{ch9:9.20}\group{chapter-9}\source{gilbarg-trudinger:page-0256}
Let $u\in W^{2,n}(\Omega)$ satisfy $Lu\ge f$ with $f\in L^n(\Omega)$.
For every $B_{2R}(y)\subset\Omega$ and $p>0$,
\begin{equation}\tag{9.48}
\sup_{B_R(y)}u\le C\left[
\left(\frac1{|B_{2R}|}\int_{B_{2R}(y)}(u^+)^p\right)^{1/p}
+\frac R\lambda\|f\|_{L^n(B_{2R}(y))}\right],
\end{equation}
where $C=C(n,\gamma,\nu R^2,p)$.
\end{theorem}
\begin{proof}
\sketch Scale to the unit ball and set
\begin{equation}
\tag{9.49}
\eta(x)=(1-|x|^2)^\beta,
\qquad v=\eta u.
\end{equation}
On the upper contact set of $v$, concavity controls
$Du$ by a multiple of $(1-|x|^2)^{-1}u$.  Lemma
\ref{gt-ch09-lem-contact-set-principal-part} then bounds $\sup v$ by an
$L^n$ norm of a power of $u^+$ and the inhomogeneous term.  Choose
$\beta=2n/p$ and use Young's inequality; larger $p$ follows from a smaller
exponent.  Approximation gives the result for $W^{2,n}$ solutions.
\end{proof}

\begin{corollary}[Local minimum principle]
\label{gt-ch09-cor-local-minimum-principle}
\dcref{ch9:9.21}\group{chapter-9}\source{gilbarg-trudinger:page-0258}
Let $u\in W^{2,n}(\Omega)$ satisfy $Lu\le f$, or $Lu=f$, with
$f\in L^n(\Omega)$.  For every $B_{2R}(y)\subset\Omega$ and $p>0$,
\begin{equation}\tag{9.50}
\sup_{B_R(y)}(-u)\le C\left[
\left(\frac1{|B_{2R}|}\int_{B_{2R}(y)}(u^-)^p\right)^{1/p}
+\frac R\lambda\|f\|_{L^n(B_{2R}(y))}\right],
\end{equation}
where $C=C(n,\gamma,\nu R^2,p)$.
\end{corollary}
\begin{proof}
\sketch Apply Theorem~\ref{gt-ch09-thm-local-maximum-principle} to $-u$.
\end{proof}

In particular, if $Lu\ge0$ and $u\ge0$ in $B_R(y)$, then
\begin{equation}\tag{9.51}
u(y)\le\frac C{R^n}\int_{B_R(y)}u,
\end{equation}
where $C=C(n,\gamma,\nu R^2)$.

\section*{9.8. H\"older and Harnack Estimates}

\begin{theorem}[Weak Harnack inequality]
\label{gt-ch09-thm-weak-harnack}
\dcref{ch9:9.22}\group{chapter-9}\source{gilbarg-trudinger:page-0258}
Let $u\in W^{2,n}(\Omega)$ satisfy $Lu\le f$, with $f\in L^n(\Omega)$, and
assume $u\ge0$ in $B_{2R}(y)\subset\Omega$.  There are $p,C>0$, depending
only on $n,\gamma,\nu R^2$, such that
\begin{equation}\tag{9.52}
\left(\frac1{|B_R|}\int_{B_R(y)}u^p\right)^{1/p}
\le C\left(\inf_{B_R(y)}u+\frac R\lambda
\|f\|_{L^n(B_{2R}(y))}\right).
\end{equation}
\end{theorem}
\begin{proof}
\sketch Normalize the ball and $\lambda$, and set
\begin{equation}
\tag{9.53}
\bar u=u+\varepsilon+\|f\|_n,\qquad
w=-\log\bar u,\qquad v=\eta w,\qquad g=f/\bar u.
\end{equation}
Take $\eta=(1-|x|^2)^\beta$. The contact-set estimate gives
\begin{equation}
\tag{9.54}
\sup_Bv\leq C(1+\|v^+\|_{n;B_\alpha}).
\end{equation}
The contact-set estimate bounds $\sup v$ unless the positive set of $v$ has
large relative measure. More precisely, if $K=K_r(z)$ and
\begin{equation}
\tag{9.55}
|K^+|\leq\theta|K|,
\end{equation}
then
\begin{equation}
\tag{9.56}
\sup_{K_{3nr}(z)}w\leq C(n,\gamma,\nu).
\end{equation}
Lemma~\ref{gt-ch09-lem-cube-level-set-growth} then yields the power decay
\[
|\{\bar u>t\}\cap K|\le
C(\inf_K\bar u/t)^\kappa.\tag{9.60}
\]
The distribution identity gives, for $p<\kappa$,
\begin{equation}
\tag{9.61}
\int_{B_\alpha}\bar u^p
\leq C(\inf_{B_\alpha}\bar u)^p.
\end{equation}
Let $\varepsilon\downarrow0$, cover the half-radius ball, and rescale.
\end{proof}

\begin{lemma}[Cube level-set growth]
\label{gt-ch09-lem-cube-level-set-growth}
\dcref{ch9:9.23}\group{chapter-9}\source{gilbarg-trudinger:page-0260}
Let $K_0$ be a cube, $w\in L^1(K_0)$, and
$\Gamma_k=\{x\in K_0:w(x)\le k\}$.  Suppose $0<\delta<1$ and $C>0$ are
such that
\begin{equation}\tag{9.57}
\sup_{K_0\cap K_{3r}(z)}(w-k)\le C
\end{equation}
whenever $K_r(z)\subset K_0$ and
\begin{equation}\tag{9.58}
|\Gamma_k\cap K_r(z)|\ge\delta|K_r(z)|.
\end{equation}
Then
\begin{equation}\tag{9.59}
\sup_{K_0}(w-k)\le
C\left(1+\frac{\log(|\Gamma_k|/|K_0|)}{\log\delta}\right).
\end{equation}
\end{lemma}
\begin{proof}
\sketch Induct on $m$ to show that
$|\Gamma_k|\ge\delta^m|K_0|$ implies
$\sup_{K_0}(w-k)\le mC$.  The cube decomposition and (9.20) enlarge the
controlled set by the factor $\delta^{-1}$ after raising the level by $C$.
Choose $m$ from the ratio $|\Gamma_k|/|K_0|$ to obtain (9.59).
\end{proof}

\begin{corollary}[Interior H\"older oscillation estimate]
\label{gt-ch09-cor-interior-holder-oscillation}
\dcref{ch9:9.24}\group{chapter-9}\source{gilbarg-trudinger:page-0262}
Let $u\in W^{2,n}(\Omega)$ solve $Lu=f$.  If
$B_{R_0}(y)\subset\Omega$ and $R\le R_0$, then
\begin{equation}\tag{9.62}
\operatorname{osc}_{B_R(y)}u\le
C\left(\frac R{R_0}\right)^\alpha
\left(\operatorname{osc}_{B_{R_0}(y)}u+kR_0\right),
\end{equation}
where $k=\|f-cu\|_{L^n(B_{R_0}(y))}$ and
$C,\alpha>0$ depend only on $n,\gamma,\nu R_0^2$.
\end{corollary}
\begin{proof}
\sketch Apply Theorem~\ref{gt-ch09-thm-weak-harnack} to the nonnegative
deficits from the supremum and infimum on concentric balls.  Their sum gives
a geometric oscillation decay with error $kR_0$; iteration yields (9.62).
\end{proof}

\begin{corollary}[Interior Harnack inequality]
\label{gt-ch09-cor-interior-harnack}
\dcref{ch9:9.25}\group{chapter-9}\source{gilbarg-trudinger:page-0262}
If $u\in W^{2,n}(\Omega)$ satisfies $Lu=0$ and $u\ge0$, then for every
$B_{2R}(y)\subset\Omega$,
\begin{equation}\tag{9.63}
\sup_{B_R(y)}u\le C\inf_{B_R(y)}u,
\end{equation}
where $C=C(n,\gamma,\nu R^2)$.
\end{corollary}
\begin{proof}
\sketch Combine the weak Harnack inequality with the local maximum principle,
using the weak-Harnack exponent in (9.48).
\end{proof}

\section*{9.9. Local Estimates at the Boundary}

\begin{theorem}[Boundary local maximum principle]
\label{gt-ch09-thm-boundary-local-maximum}
\dcref{ch9:9.26}\group{chapter-9}\source{gilbarg-trudinger:page-0262}
Let $u\in W^{2,n}(\Omega)\cap C^0(\overline\Omega)$ satisfy $Lu\ge f$ in
$\Omega$ and $u\le0$ on $B_{2R}(y)\cap\partial\Omega$, where
$f\in L^n(\Omega)$.  For every $p>0$,
\begin{equation}\tag{9.64}
\sup_{\Omega\cap B_R(y)}u\le C\left[
\left(\frac1{|B_{2R}|}\int_{\Omega\cap B_{2R}(y)}(u^+)^p\right)^{1/p}
+\frac R\lambda\|f\|_{L^n(\Omega\cap B_{2R}(y))}\right],
\end{equation}
where $C=C(n,\gamma,\nu R^2,p)$.
\end{theorem}
\begin{proof}
\sketch Extend $u$ by zero to the whole ball.  Although the extension need
not be $C^2$, the upper contact set of the cutoff product lies inside
$\Omega$, where the original function has the needed regularity.  The proof
of Theorem~\ref{gt-ch09-thm-local-maximum-principle} therefore applies
unchanged; approximation treats strong solutions.
\end{proof}

The same observation allows the contact set in (9.11) to be taken relative to
any larger domain to which $u$ is extended by zero, with the diameter of that
larger domain.

\begin{theorem}[Boundary weak Harnack inequality]
\label{gt-ch09-thm-boundary-weak-harnack}
\dcref{ch9:9.27}\group{chapter-9}\source{gilbarg-trudinger:page-0263}
Let $u\in W^{2,n}(\Omega)$ satisfy $Lu\le f$ in $\Omega$ and $u\ge0$ in
$B_{2R}(y)\cap\Omega$.  Set
\[
m=\inf_{B_{2R}(y)\cap\partial\Omega}u
\]
and extend the truncation by
\[
u_m^-(x)=
\begin{cases}
\min\{u(x),m\},&x\in B_{2R}(y)\cap\Omega,\\
m,&x\in B_{2R}(y)\setminus\Omega.
\end{cases}
\]
Then, for constants $p,C>0$ depending only on $n,\gamma,\nu R^2$,
\begin{equation}\tag{9.65}
\left(\frac1{|B_R|}\int_{B_R(y)}(u_m^-)^p\right)^{1/p}
\le C\left(\inf_{\Omega\cap B_R(y)}u+
\frac R\lambda\|f\|_{L^n(\Omega\cap B_{2R}(y))}\right).
\end{equation}
If only $u\in W^{2,n}_{\mathrm{loc}}(\Omega)$ is assumed, replace $m$ by
$\liminf_{x\to B_{2R}(y)\cap\partial\Omega}u(x)$.
\end{theorem}
\begin{proof}
\sketch Repeat the proof of Theorem~\ref{gt-ch09-thm-weak-harnack} with the
extended truncation.  The level-set estimate applies up to level $m$, while
the distribution function vanishes above $m$.  The liminf formulation
follows by interior exhaustion.
\end{proof}

\begin{corollary}[Boundary oscillation estimate under an exterior cone]
\label{gt-ch09-cor-boundary-oscillation-exterior-cone}
\dcref{ch9:9.28}\group{chapter-9}\source{gilbarg-trudinger:page-0263}
Let $u\in W^{2,n}_{\mathrm{loc}}(\Omega)$ solve $Lu=f$, with
$f\in L^n(\Omega)$, and suppose $\Omega$ has an exterior cone $V_y$ at
$y\in\partial\Omega$.  For $0<R<R_0$ and $B_0=B_{R_0}(y)$,
\begin{equation}\tag{9.66}
\operatorname{osc}_{\Omega\cap B_R(y)}u
\le C\left[
\left(\frac R{R_0}\right)^\alpha
\left(\operatorname{osc}_{\Omega\cap B_0}u+kR_0\right)
+\sigma(\sqrt{RR_0})\right],
\end{equation}
where
\[
k=\|f-cu\|_{L^n(\Omega\cap B_0)}
\]
and
\[
\sigma(r)=\operatorname{osc}_{\partial\Omega\cap B_r(y)}u
=\limsup_{x\to\partial\Omega\cap B_r(y)}u(x)
-\liminf_{x\to\partial\Omega\cap B_r(y)}u(x).
\]
The positive constants $C,\alpha$ depend only on
$n,\gamma,\nu R_0^2,V_y$.
\end{corollary}
\begin{proof}
\sketch The exterior cone supplies a fixed positive proportion of every
small boundary ball on which the extended truncations are constant.  Apply
Theorem~\ref{gt-ch09-thm-boundary-weak-harnack} to upper and lower deficits,
then iterate the resulting oscillation decay from $R_0$ to $R$.  Splitting
at the geometric-mean scale produces the boundary term
$\sigma(\sqrt{RR_0})$.
\end{proof}

\begin{corollary}[Global H\"older estimate under a uniform exterior cone]
\label{gt-ch09-cor-global-holder-exterior-cone}
\dcref{ch9:9.29}\group{chapter-9}\source{gilbarg-trudinger:page-0264}
Let $u\in W^{2,n}(\Omega)\cap C^0(\overline\Omega)$ solve $Lu=f$, where
$f\in L^n(\Omega)$, and let $u=\varphi$ on $\partial\Omega$ with
$\varphi\in C^\beta(\overline\Omega)$ for some $\beta>0$.  If
$\partial\Omega$ satisfies a uniform exterior cone condition, then
$u\in C^\alpha(\overline\Omega)$ and
\begin{equation}\tag{9.67}
|u|_{\alpha;\Omega}\le C,
\end{equation}
where $C,\alpha>0$ depend on
$n,\gamma,\nu,\beta,\Omega,|\varphi|_{\beta;\Omega}$, and
$|u|_{0;\Omega}$.
\end{corollary}
\begin{proof}
\sketch Insert the boundary H\"older modulus of $\varphi$ into Corollary
\ref{gt-ch09-cor-boundary-oscillation-exterior-cone}; the uniform cone gives
uniform boundary constants.  Combine the resulting estimates with Corollary
\ref{gt-ch09-cor-interior-holder-oscillation} on an interior covering.
\end{proof}

\begin{theorem}[Dirichlet problem under an exterior cone condition]
\label{gt-ch09-thm-dirichlet-exterior-cone}
\dcref{ch9:9.30}\group{chapter-9}\source{gilbarg-trudinger:page-0264}
Let $\Omega$ be bounded and satisfy an exterior cone condition at every
boundary point.  Suppose $L$ is strictly elliptic,
$a^{ij}\in C^0(\Omega)\cap L^\infty(\Omega)$,
$b^i,c\in L^\infty(\Omega)$, and $c\le0$.  If $p\ge n$,
$f\in L^p(\Omega)$, and $\varphi\in C^0(\partial\Omega)$, then
\[
Lu=f\quad\text{in }\Omega,\qquad u=\varphi\quad\text{on }\partial\Omega
\]
has a unique solution
$u\in W^{2,p}_{\mathrm{loc}}(\Omega)\cap C^0(\overline\Omega)$.
\end{theorem}
\begin{proof}
\sketch Construct a Perron family using local solutions in balls from
Corollary~\ref{gt-ch09-cor-strong-dirichlet-continuous-data}.
The strong maximum principle supplies comparison, Theorem
\ref{gt-ch09-thm-interior-w2p-estimate} and weak compactness give a local
$W^{2,p}$ Perron solution, and Corollary
\ref{gt-ch09-cor-boundary-oscillation-exterior-cone} enforces the continuous
boundary values.  Theorem~\ref{gt-ch09-thm-strong-dirichlet-uniqueness}
gives uniqueness.
\end{proof}

\subsection*{Boundary H\"older Estimates for the Gradient}

Assume here that $Lu=a^{ij}D_{ij}u$ satisfies (9.47).

\begin{theorem}[Boundary H\"older estimate for the normal quotient]
\label{gt-ch09-thm-boundary-holder-normal-quotient}
\dcref{ch9:9.31}\group{chapter-9}\source{gilbarg-trudinger:page-0265}
Let $B^+=B_{R_0}(0)\cap\mathbb R^n_+$ and
$T=B_{R_0}(0)\cap\partial\mathbb R^n_+$.  Suppose
$u\in W^{2,n}_{\mathrm{loc}}(B^+)\cap C^0(\overline{B^+})$ satisfies
$Lu=f$ with $f\in L^\infty(B^+)$ and $u=0$ on $T$.  For $R\le R_0$,
\begin{equation}\tag{9.68}
\operatorname{osc}_{B_R^+}\frac u{x_n}\le
C\left(\frac R{R_0}\right)^\alpha
\left(\operatorname{osc}_{B^+}\frac u{x_n}
+R_0\sup_{B^+}|f/\lambda|\right),
\end{equation}
where $C,\alpha>0$ depend only on $n,\gamma$.  The boundary gradient exists,
is H\"older continuous on $T$, and satisfies
\begin{equation}\tag{9.71}
\operatorname{osc}_{|x'|<R}Du(x',0)\le
C\left(\frac R{R_0}\right)^\alpha
\left(\operatorname{osc}_{B^+}\frac u{x_n}
+\frac R\lambda\sup_{B^+}|f|\right).
\end{equation}
In either estimate, $\operatorname{osc}_{B^+}(u/x_n)$ may be replaced by
$\operatorname{osc}_{B^+}u/R_0$ or by $\sup_{B^+}|Du|$.
\end{theorem}
\begin{proof}
\sketch Write $v=u/x_n$.  A thin-slab barrier and the maximum principle give
\begin{equation}
\tag{9.69}
\inf_{\substack{|x'|<R\\x_n=\delta R}}v
\leq2\left(\inf_{B_{R/2,\delta}}v
+\frac R\lambda\sup_{B^+}|f|\right),
\end{equation}
where $B_{R,\delta}=\{x:|x'|<R,\ 0<x_n<\delta R\}$ and
$\delta=\delta(n,\gamma)>0$. Apply
Corollary~\ref{gt-ch09-cor-interior-harnack} in a
strip bounded away from $x_n=0$ to get
\[
\sup v\le C\left(\inf v+R\sup|f/\lambda|\right).\tag{9.70}
\]
Apply this to the upper and lower deficits of $v$ to obtain geometric
oscillation decay, then iterate.  The trace of $v$ is the normal derivative,
while tangential derivatives vanish, giving (9.71).
\end{proof}
